\scanpage{83}
de $\unsure{\kappa|k}$). On peut dire aussi que la donnée de $\kappa$ équivaut à celle d'une \deleted{paire} \inserted{couple} $(\alpha.\,\kappa)$, où $\alpha$ est une section de $\mathord{\mathrm{I}\mkern-7mu\mu}_4^*\wedge\underline\omega_{\underline J}$, et où $\kappa$ est un iso. (31) \uline{compatible} avec $\alpha$ au sens évident.

Pour $\alpha$ fixé, les $\kappa$ correspondants forment les sections d'un torseur sous $\mathord{\mathrm{I}\mkern-7mu\mu}_2$, -- ou \uline{torseur des épinglages de Legendre} de $E/S$, \deleted{\unsure{muni de|munis de}} \inserted{associés à} ``l'épinglage partiel'' $\alpha$.

La catégorie \inserted{\deleted{groupoïde}} \deleted{\unsure{de}}fibrée (qui est un champ en groupoïdes) des courbes elliptiques relatives, munies d'un épinglage de Legendre, équivaut à celle des triplets $(\underline J,t,\alpha)$, où

\[
\left\{\begin{aligned}
&\text{a) }\underline J\quad\text{rev. étale cubique de }S\\
&\text{b) }t\ \text{section de }\Sigma_{\underline J}^*\\
&\text{c) }\alpha\ \text{section de }\mathord{\mathrm{I}\mkern-7mu\mu}_4^*\wedge\underline\omega_{\underline J},\ \text{i.e iso. }\mu_4^*\simeq\underline\omega_{\underline J}
\end{aligned}\right.
\tag*{(\unsure{34|39})}
\]

% slanted
\margin{\uline{NB} Il faudrait expliciter la construction de $E$, en termes de ces données\ldots{} (via fonctions $\wp$ de \unsure{Weierstrass}}
% archive sheet number: 82
